بگذارید
$$
N=2023^{2023}.
$$

ادعا می‌کنیم که چندجمله‌ای‌های مطلوب دقیقاً چندجمله‌ای‌های زیر هستند:

\begin{itemize}
\item $P(x)=x^m$ برای یک $m\in\mathbb{N}$؛
\item $P(x)=c$ برای یک $c\in\mathbb{N}$ که $\omega(c)\leq N+1$ باشد.
\end{itemize}

ابتدا بررسی می‌کنیم که هر یک از این چندجمله‌ای‌ها شرط مسئله را ارضا می‌کند. اگر $P(x)=x^m$ باشد، آنگاه
$$
\omega(P(n))=\omega(n).
$$
همچنین اگر $P(x)=c$ و $\omega(c)\leq N+1$ باشد، هرگاه $\omega(n)>N$ داشته باشیم،
$$
\omega(n)\geq N+1\geq\omega(c).
$$

اکنون ثابت می‌کنیم که چندجمله‌ای دیگری وجود ندارد.

اگر $P(x)=c$ ثابت باشد، چون $P(n)$ باید برای اعداد طبیعی موردنظر مثبت باشد، داریم $c\in\mathbb{N}$. همچنین اگر $\omega(c)>N+1$ باشد، با انتخاب $n$ به صورت حاصل‌ضرب $N+1$ عدد اول متمایز، خواهیم داشت
$$
\omega(n)=N+1<\omega(c),
$$
که با شرط مسئله در تناقض است. پس
$$
\omega(c)\leq N+1.
$$

اکنون فرض کنید $P$ ثابت نباشد. می‌توانیم بنویسیم
$$
P(x)=\alpha x^mQ(x),
$$
که در آن $\alpha\in\mathbb{Z}\setminus\{0\}$، $m\geq0$ و
$$
Q(0)\neq0.
$$

ابتدا لم زیر را ثابت می‌کنیم.

\textbf{لم.}
اگر $f(x)\in\mathbb{Z}[x]$ چندجمله‌ای غیرثابت باشد، آنگاه بی‌نهایت عدد اول $q$ وجود دارد که برای یک عدد صحیح $a$ داشته باشیم
$$
q\mid f(a).
$$

\textbf{اثبات.}
اگر $f(0)=0$ باشد، نتیجه بدیهی است؛ زیرا هر عدد اولی که یک مقدار ناصفر از $f$ را تقسیم کند، مقسوم‌علیه اول یکی از مقادیر $f$ است.

حال فرض کنید $f(0)\neq0$. چندجمله‌ای
$$
g(x)=\frac{f(xf(0))}{f(0)}
$$
را در نظر بگیرید. چون جملهٔ ثابت $f(xf(0))$ برابر $f(0)$ است، چندجمله‌ای $g(x)$ ضرایب صحیح دارد و
$$
g(0)=1.
$$

فرض کنید تنها تعداد متناهی عدد اول $p_1,\ldots,p_k$ وجود داشته باشد که مقسوم‌علیه یکی از مقادیر $g$ باشند. عدد صحیح $M$ را آن‌قدر بزرگ انتخاب می‌کنیم که
$$
|g(Mp_1\cdots p_k)|>1.
$$
از طرفی برای هر $i$ داریم
$$
g(Mp_1\cdots p_k)\equiv g(0)=1\pmod{p_i},
$$
پس هیچ‌یک از $p_i$ها مقدار $g(Mp_1\cdots p_k)$ را تقسیم نمی‌کند. اما این عدد قدرمطلقی بزرگ‌تر از $1$ دارد و بنابراین دارای یک مقسوم‌علیه اول است؛ تناقض.

پس بی‌نهایت عدد اول مقادیر $g$ و در نتیجه بی‌نهایت عدد اول مقادیر $f$ را تقسیم می‌کنند.
\hfill$\square$

اکنون فرض کنید $Q$ ثابت باشد، مثلاً $Q(x)=c$. در این صورت
$$
P(x)=\alpha cx^m.
$$
اگر $d=\alpha c$ بگذاریم، داریم
$$
P(x)=dx^m.
$$

از آنجا که $P(n)$ باید برای تمام اعداد طبیعی $n$ با $\omega(n)>N$ مثبت باشد، باید $d>0$ باشد.

اگر $d>1$ باشد، $n$ را حاصل‌ضرب $N+1$ عدد اول متمایز انتخاب می‌کنیم که هیچ‌کدام با $d$ مقسوم‌علیه مشترک نداشته باشند. در این صورت
$$
\omega(n)=N+1,
$$
ولی
$$
\omega(P(n))=\omega(dn^m)=\omega(d)+\omega(n)>N+1,
$$
که با شرط مسئله تناقض دارد. بنابراین $d=1$ و در نتیجه
$$
P(x)=x^m.
$$

اکنون تنها حالت باقی‌مانده این است که $Q$ غیرثابت باشد. بنویسید
$$
Q(x)=\alpha_dx^d+\cdots+\alpha_1x+\alpha_0,
$$
که در آن
$$
\alpha_d\neq0,\qquad\alpha_0\neq0.
$$

طبق لم، بی‌نهایت عدد اول وجود دارد که مقسوم‌علیه یکی از مقادیر $Q$ هستند. بنابراین می‌توانیم $T>N+1$ عدد اول متمایز
$$
q_1,\ldots,q_T
$$
را طوری انتخاب کنیم که
$$
q_i>|\alpha_0|
$$
برای هر $i$.

برای هر $i$ عدد صحیح ناصفر $a_i$ای انتخاب می‌کنیم که
$$
q_i\mid Q(a_i).
$$
می‌توان $a_i$ را ناصفر انتخاب کرد؛ زیرا اگر $q_i\mid a_i$ باشد، آنگاه
$$
q_i\mid Q(a_i)-Q(0),
$$
و چون $q_i\mid Q(a_i)$، نتیجه می‌شود
$$
q_i\mid Q(0)=\alpha_0,
$$
که با $q_i>|\alpha_0|$ تناقض دارد.

اکنون $N+1$ عدد اول متمایز
$$
p_1,\ldots,p_{N+1}
$$
را طوری انتخاب می‌کنیم که
$$
p_1\equiv a_i\pmod{q_i}
\qquad(1\leq i\leq T)
$$
و
$$
p_j\equiv1\pmod{q_1q_2\cdots q_T}
\qquad(2\leq j\leq N+1).
$$

وجود چنین اعداد اولی از قضیهٔ باقی‌مانده‌های چینی و قضیهٔ دیریکله دربارهٔ اعداد اول در تصاعدهای حسابی نتیجه می‌شود.

حال قرار دهید
$$
n=p_1p_2\cdots p_{N+1}.
$$
در این صورت
$$
\omega(n)=N+1.
$$

از طرفی برای هر $i$ داریم
$$
n\equiv p_1\equiv a_i\pmod{q_i},
$$
پس
$$
Q(n)\equiv Q(a_i)\equiv0\pmod{q_i}.
$$
بنابراین هر یک از $q_i$ها مقسوم‌علیه $Q(n)$ و در نتیجه مقسوم‌علیه $P(n)$ است. پس
$$
\omega(P(n))\geq T.
$$
اما $T>N+1$ است، بنابراین
$$
\omega(P(n))>N+1=\omega(n),
$$
که با شرط مسئله تناقض دارد.

پس $Q$ حتماً ثابت است و در نتیجه تنها چندجمله‌ای‌های غیرثابت مطلوب از فرم
$$
P(x)=x^m,\qquad m\in\mathbb{N}
$$
هستند.

در نتیجه تمام جواب‌ها عبارت‌اند از
$$
\boxed{P(x)=x^m\quad(m\in\mathbb{N})}
$$
و
$$
\boxed{P(x)=c\quad(c\in\mathbb{N},\ \omega(c)\leq2023^{2023}+1)}.
$$